\section{CNN 的广泛应用}

\subsection{超越分类：一种通用架构}

CNN 的特征学习部分是高度通用的——学到的特征表示可以用于多种不同的下游任务：

\begin{figure}[H]
\centering
\includegraphics[width=0.85\textwidth]{slides-images/slide-56.jpg}
\caption{CNN 作为通用架构：同样的特征学习模块可以支持分类、目标检测、语义分割、概率控制等多种任务。\protect\footnotemark}
\end{figure}
\footnotetext{Slides 第56页。}

\begin{itemize}
    \item \textbf{分类}（Classification）：判断整张图像属于哪个类别
    \item \textbf{目标检测}（Object Detection）：定位并分类图像中的所有物体
    \item \textbf{语义分割}（Segmentation）：对每个像素进行分类
    \item \textbf{概率控制}（Probabilistic Control）：如自动驾驶中从图像到方向盘转角
\end{itemize}

\subsection{应用一：医学影像中的乳腺癌筛查}

CNN 在医学影像分析中取得了突破性成果。Amini 展示了一项发表在 Nature 上的研究，证明基于 CNN 的 AI 系统在乳腺癌筛查中\textbf{超越了专业放射科医生}的表现：

\begin{figure}[H]
\centering
\includegraphics[width=0.9\textwidth]{slides-images/slide-57.jpg}
\caption{CNN 用于乳腺癌筛查：ROC 曲线显示 AI 系统的灵敏度-特异度均优于放射科医生。右图展示了一例人类医生遗漏但 AI 检出的乳腺癌病灶。\protect\footnotemark}
\end{figure}
\footnotetext{Slides 第57页。参考 McKinney+ Nature 2020。}

\begin{importantbox}{CNN 在医学影像中的价值}
\begin{itemize}
    \item 在美国乳腺癌 2 年筛查和 1 年筛查的评估中，AI 系统的 ROC 曲线均优于放射科医生群体
    \item AI 能检出人类医生遗漏的微小病灶，具有重要的临床辅助价值
    \item 类似的 CNN 方法已扩展到皮肤癌检测（Esteva+ Nature 2017）和 COVID-19 影像诊断
\end{itemize}
\end{importantbox}

\subsection{应用二：目标检测}

\textbf{目标检测}（Object Detection）比分类更进一步——不仅要识别物体类别，还要给出物体的精确位置（bounding box）。

\begin{figure}[H]
\centering
\includegraphics[width=0.85\textwidth]{slides-images/slide-58.jpg}
\caption{分类 vs 目标检测：分类只输出类别标签（如 Taxi）；目标检测同时输出类别和边界框位置 $(x, y, w, h)$。\protect\footnotemark}
\end{figure}
\footnotetext{Slides 第58页。}

目标检测的输出格式为每个检测到的物体一组 $(class, x, y, w, h)$：
\begin{itemize}
    \item $class$：物体类别
    \item $(x, y)$：边界框左上角坐标
    \item $(w, h)$：边界框的宽度和高度
\end{itemize}

\begin{figure}[H]
\centering
\includegraphics[width=0.85\textwidth]{slides-images/slide-59.jpg}
\caption{多目标检测：模型需要输出场景中\textbf{每个}物体的类别和边界框，且物体数量是不固定的。\protect\footnotemark}
\end{figure}
\footnotetext{Slides 第59页。}

\begin{warningbox}{目标检测的难点}
目标检测比分类困难得多，主要原因是：
\begin{enumerate}
    \item 边界框可以出现在图像的\textbf{任何位置}
    \item 边界框可以是\textbf{任意大小}
    \item 场景中物体的\textbf{数量不固定}——可以是零个、一个或多个
    \item 不同物体可能有不同的类别
\end{enumerate}
这使得输出空间的维度不再固定，需要特殊的网络设计来处理。
\end{warningbox}

\subsubsection{朴素方法：滑动窗口}

最直觉的方法是在图像上尝试所有可能的位置和大小的边界框，对每个框内的图像区域用 CNN 分类。但这种方法的计算量是灾难性的——可能的窗口数量随位置、大小和宽高比呈指数增长。

\begin{figure}[H]
\centering
\includegraphics[width=0.85\textwidth]{slides-images/slide-60.jpg}
\caption{朴素目标检测：遍历所有可能的窗口，每个送入 CNN 分类。问题：窗口数量太多，不同位置、尺度、大小组合产生海量候选区域。\protect\footnotemark}
\end{figure}
\footnotetext{Slides 第60页。}

\subsubsection{R-CNN 系列}

为了解决效率问题，研究者提出了 R-CNN（Region-based CNN）系列方法。核心思想是先用选择性搜索（Selective Search）等方法生成少量候选区域（region proposals），然后只对这些候选区域进行 CNN 分类。后续发展包括 Fast R-CNN 和 Faster R-CNN，进一步将区域提议网络（Region Proposal Network, RPN）集成到 CNN 中，实现端到端训练。

\begin{knowledgebox}{R-CNN 到 Faster R-CNN 的演进}
\begin{itemize}
    \item \textbf{R-CNN}（2014）：选择性搜索生成约 2000 个候选区域，每个独立通过 CNN 提取特征，再用 SVM 分类。速度很慢。
    \item \textbf{Fast R-CNN}（2015）：先对整幅图像做一次 CNN 特征提取，再从特征图上裁剪候选区域。共享了特征计算，速度大幅提升。
    \item \textbf{Faster R-CNN}（2015）：引入 Region Proposal Network (RPN) 直接在特征图上生成候选区域，完全端到端训练。
\end{itemize}
\end{knowledgebox}

\subsection{本章小结}

CNN 的特征学习模块是通用的，可以支持分类、目标检测、分割等多种视觉任务。在医学影像中，CNN 已经达到甚至超越人类专家水平。目标检测需要同时预测类别和位置，朴素的滑动窗口方法效率低下，R-CNN 系列通过区域提议机制大幅提升了效率。

\newpage
